% !TEX root
% !TEX program = xelatex
% !BIB program = biber

% \def \PrintMode{} %在使用电子版论文时，请将此行注释。在打印纸质论文时，请保持本行命令不被注释，然后打印时选择双面打印即可。

%用来控制是否启动打印模式的宏，请勿改动。
\ifx \PrintMode \undefined
    \def \SideMode{oneside}
    \def \ClearPageStyle{\clearpage}
\else
    \def \SideMode{twoside}
    \def \ClearPageStyle{\cleardoublepage}
\fi

\documentclass[a4paper,\SideMode,UTF8]{article} %A4纸，UTF-8

\usepackage[thmmarks,hyperref]{ntheorem} %定义命令环境使用的宏包
\usepackage[heading,zihao=-4]{ctex} %用来提供中文支持
\usepackage{amsmath,amssymb} %数学符号等相关宏包
\usepackage{graphicx} %插入图片所需宏包
\usepackage{xspace} %提供一些好用的空格命令
\usepackage{tikz-cd} %画交换图需要的宏包
\usepackage{url} %更好的超链接显示
\usepackage{array,booktabs} %表格相关的宏包
\usepackage{caption} %实现图片的多行说明
\usepackage{float} %图片与表格的更好排版
\usepackage{ulem} %更好的下划线
\usepackage{multirow}

\usepackage[top=2.5cm, bottom=2.0cm, left=3.0cm, right=2.0cm]{geometry} %设置页边距

\usepackage{fontspec} %设置字体需要的宏包

%设置西文字体为Times New Roman，如果没有则以开源近似字体代替
\IfFontExistsTF{Times New Roman}{
	\setmainfont{Times New Roman}
}{
	\usepackage{newtxtext,newtxmath}
}

%设置文档中文字体。优先次序：中易 > Adobe > 华文(Mac) > Fandol
\IfFontExistsTF{SimSun}{
	\setCJKmainfont[AutoFakeBold=2,ItalicFont=KaiTi]{SimSun}
}{
	\IfFontExistsTF{AdobeSongStd-Light}{
		\setCJKmainfont[AutoFakeBold=2,ItalicFont=AdobeKaitiStd-Regular]{AdobeSongStd-Light}
	}{
		\IfFontExistsTF{STSong}{
			\setCJKmainfont[AutoFakeBold=2,BoldFont=STHeiti,ItalicFont=STKaiti]{STSong}
		}{
			\setCJKmainfont[AutoFakeBold=2,ItalicFont=FandolKai-Regular]{FandolSong-Regular}
		}
	}
}
\IfFontExistsTF{SimHei}{
	\setCJKsansfont[AutoFakeBold=2]{SimHei}
}{
	\IfFontExistsTF{AdobeHeitiStd-Regular}{
		\setCJKsansfont[AutoFakeBold=2]{AdobeHeitiStd-Regular}
	}{
		\IfFontExistsTF{STHeiti}{
			\setCJKsansfont [AutoFakeBold=2]{STHeiti}
		}{
			\setCJKsansfont[AutoFakeBold=2]{FandolHei-Regular}
		}
	}
}


\IfFileExists{zhlineskip.sty}{
	%Microsoft Word 样式的1.5倍行距（按中易字体计算）
	\usepackage[
		restoremathleading=false,
		UseMSWordMultipleLineSpacing,
		MSWordLineSpacingMultiple=1.5
	]{zhlineskip}
}{
	\linespread{1.621} %1.5倍行距
}

\showboxdepth=5
\showboxbreadth=5

%设置各级系统的编号格式
\setcounter{secnumdepth}{5}
\ctexset { section = { name={,.},number={\arabic{section}},format={\sffamily \zihao {-4}} } }
\ctexset { subsection = { name={,},number={\arabic{section}.\arabic{subsection}},format={\sffamily \zihao {-4}} } }
\ctexset { subsubsection = { name={,},number={\arabic{section}.\arabic{subsection}.\arabic{subsubsection}},format={\sffamily \zihao {-4}},indent=2em } }
\ctexset { paragraph = { name={,},number={\arabic{section}.\arabic{subsection}.\arabic{subsubsection}.\arabic{paragraph}},format={\sffamily \zihao {-4}},indent=4em } }
\ctexset { subparagraph = { name={,)},number={\arabic{subparagraph}},format={\sffamily \zihao {-4}},indent=6em } }

\usepackage[bottom,perpage]{footmisc}               %脚注，显示在每页底部，编号按页重置
\renewcommand*{\footnotelayout}{\zihao{-5}\rmfamily}  %设置脚注为小五号宋体
\renewcommand{\thefootnote}{\arabic{footnote}}    %设置脚注标记为①，②，...

%设置页眉页脚
\usepackage{fancyhdr}

% \lhead{华东师范大学学士学位论文}
\lhead{}
\chead{华东师范大学2017级本科毕业论文/设计}
\rhead{}
% \rhead{\TitleCHS}

\lfoot{}
\cfoot{\thepage}
\rfoot{}

\usepackage{xcolor} %彩色的文字

\usepackage[hidelinks]{hyperref} %各种超链接必备
\usepackage{cleveref} %交叉引用

%设置尾注
\usepackage{endnotes}
\renewcommand{\enotesize}{\zihao{-5}}
\renewcommand{\notesname}{\sffamily \zihao {-4} 尾注}
\renewcommand\enoteformat{
	\raggedright
	\leftskip=1.8em
	\makebox[0pt][r]{\theenmark. \rule{0pt}{\dimexpr\ht\strutbox+\baselineskip}}
}
\renewcommand\makeenmark{\textsuperscript{[尾注\theenmark]}}
% \renewcommand{\thempfootnote}{\arabic{mpfootnote}}
% \usepackage{footnotebackref}

%定义证明与解环境
\theoremstyle{nonumberplain}
\theorembodyfont{\upshape}
\theoremseparator{:}
\theoremsymbol{\ensuremath{\square}}
\newtheorem{proof}{\bfseries \sffamily 证明}
\theoremsymbol{\ensuremath{\blacksquare}}
\newtheorem{solution}{\bfseries \sffamily 解}

%定义各种常用环境
\theoremstyle{plain}
\theoremseparator{.}
\theorembodyfont{\upshape}
\theoremsymbol{}
\newtheorem{theorem}{\bfseries \sffamily 定理}[section]
\renewtheorem*{theorem*}{\bfseries \sffamily 定理}
\newtheorem{lemma}[theorem]{\bfseries \sffamily 引理}
\renewtheorem*{lemma*}{\bfseries \sffamily 引理}
\newtheorem{corollary}[theorem]{\bfseries \sffamily 推论}
\renewtheorem*{corollary*}{\bfseries \sffamily 推论}
\newtheorem{definition}[theorem]{\bfseries \sffamily 定义}
\renewtheorem*{definition*}{\bfseries \sffamily 定义}
\newtheorem{conjecture}[theorem]{\bfseries \sffamily 猜想}
\renewtheorem*{conjecture*}{\bfseries \sffamily 猜想}
\newtheorem{problem}[theorem]{\bfseries \sffamily 问题}
\renewtheorem*{problem*}{\bfseries \sffamily 问题}
\newtheorem{proposition}[theorem]{\bfseries \sffamily 命题}
\renewtheorem*{proposition*}{\bfseries \sffamily 命题}
\newtheorem{remark}[theorem]{\bfseries \sffamily 注记}
\renewtheorem*{remark*}{\bfseries \sffamily 注记}
\newtheorem{example}[theorem]{\bfseries \sffamily 例}
\renewtheorem*{example*}{\bfseries \sffamily 例}

%设置各种常用环境的交叉引用格式
\crefformat{theorem}{#2\bfseries{\sffamily 定理} #1#3}
\crefformat{lemma}{#2\bfseries{\sffamily 引理} #1#3}
\crefformat{corollary}{#2\bfseries{\sffamily 推论} #1#3}
\crefformat{definition}{#2\bfseries{\sffamily 定义} #1#3}
\crefformat{conjecture}{#2\bfseries{\sffamily 猜想} #1#3}
\crefformat{problem}{#2\bfseries{\sffamily 问题} #1#3}
\crefformat{proposition}{#2\bfseries{\sffamily 命题} #1#3}
\crefformat{remark}{#2\bfseries{\sffamily 注记} #1#3}
\crefformat{example}{#2\bfseries{\sffamily 例} #1#3}

%允许公式跨页显示
\allowdisplaybreaks

%屏蔽无关的Warning
\usepackage{silence}
\WarningFilter*{biblatex}{Conflicting options.\MessageBreak'eventdate=iso' requires 'seconds=true'.\MessageBreak Setting 'seconds=true'}

%使用biblatex管理文献，输出格式使用gb7714-2015标准，后端为biber
\usepackage[backend=biber,gbpub=false,style=gb7714-2015,hyperref=true]{biblatex}
%将参考文献字体设置为五号
\renewcommand*{\bibfont}{\zihao{5}}

%生成感谢，请勿改动
\newcommand{\makeacknowledgement}{
	\clearpage
	\input{./ending/acknowledgement.tex}
}

%For Algorithm
\usepackage{algorithm,algorithmicx,algpseudocode}
\floatname{algorithm}{算法}
\renewcommand{\algorithmicrequire}{\textbf{输入:}}
\renewcommand{\algorithmicensure}{\textbf{输出:}}

%可能会需要在用自然语言描述算法步骤时使用的宏包
\usepackage{enumitem}

%表格单元格内换行
\newcommand{\tabincell}[2]{\begin{tabular}{@{}#1@{}}#2\end{tabular}}

%设置图、表的编号格式
\renewcommand{\thefigure}{\arabic{section}-\arabic{figure}}
\renewcommand{\thetable}{\arabic{section}-\arabic{table}}
%%每个section开始重置图、表的计数器
\makeatletter
\@addtoreset{table}{section}
\makeatother
\makeatletter
\@addtoreset{figure}{section}
\makeatother

%显示 1、2级标题
\setcounter{tocdepth}{2}

%设置目录字体
\usepackage{tocloft}
\renewcommand{\contentsname}{\centerline{目录}}
\renewcommand{\cftaftertoctitle}{\hfill}
\renewcommand{\cfttoctitlefont}{\sffamily \bfseries \zihao{-3}}
\renewcommand{\cftsubsubsecfont}{\rmfamily}
\renewcommand{\cftsubsecfont}{\rmfamily}
\renewcommand{\cftsecfont}{\rmfamily}
\renewcommand{\cftsecleader}{\cftdotfill{\cftdotsep}}
\renewcommand{\cftsecfont}{}
\renewcommand{\cftsecpagefont}{}

%灵活的行距定义（用于封面）
\usepackage{setspace}
%使用绝对坐标制作封面使用的宏包
\usepackage[absolute,overlay]{textpos}
  \setlength{\TPHorizModule}{1mm}
  \setlength{\TPVertModule}{1mm}
 %加载各宏包以及本模板的主要设置
\addbibresource{./reference/thesis-ref.bib} %加载bib文件（参考文献）

\begin{document}

\pagestyle{empty} %不对正文前的各页面使用页眉页脚
\newgeometry{top=2.0cm, bottom=2.0cm,left=3.18cm, right=3.18cm} %设置用于首页的页边距

%请不要修改本页的任何代码！
%请不要修改本页的任何代码！
%请不要修改本页的任何代码！
\thispagestyle{empty}
\begin{titlepage}
	\captionsetup{belowskip=0pt}
	\newcommand{\TitleCHS}{ 中文名字 } %中文标题

\newcommand{\TitleENG}{ English name } %英文标题

\newcommand{\Author}{xxx} %作者名字

\newcommand{\StudentID}{xxx} %学号

\newcommand{\Department}{数据科学与工程学院} %学院

\newcommand{\Major}{数据科学与大数据技术} %专业

\newcommand{\Supervisor}{xxx} %导师名字

\newcommand{\AcademicTitle}{xxx} %导师职称

\newcommand{\CompleteYear}{2021} %毕业年份

\newcommand{\CompleteMonth}{7} %毕业月份

\newcommand{\KeywordsCHS}{ 深度学习，xxx，xxx，xxx，分布式训练 } %中文关键词

\newcommand{\KeywordsENG}{ Deep Learning Framework, xxx, xxx, xxx, Distributed Training } %英文关键词

	\renewcommand{\ULthickness}{1.2pt}
	\begin{center}\noindent \bfseries \zihao{4}{\rmfamily{\CompleteYear 届本科生学士学位论文\hfill 学校代码：\uline{10269}}}\end{center}

	\begin{textblock}{146.4}(31.8,31)
		\centering
		\includegraphics{./figures/inner-cover(contains_font).eps}
	\end{textblock}

	\begin{textblock}{146.4}(31.8,125)
		\setstretch{2.0}
		\noindent
		\begin{minipage}[t][8.2cm][c]{\linewidth}
			\begin{center}
				\noindent\textbf{\zihao{1}{\rmfamily{\expandafter\uline\expandafter{\TitleCHS}}}}
			\end{center}
			\begin{center}
				\noindent\textbf{\zihao{1}{\rmfamily{\expandafter\uline\expandafter{\TitleENG}}}}
			\end{center}
		\end{minipage}
	\end{textblock}

	\renewcommand{\ULthickness}{0.4pt}

	\begin{textblock}{146.4}(31.8,208)
		\begin{center}
			\renewcommand{\arraystretch}{0.9}
			\bfseries\zihao{4}\rmfamily
			\begin{tabular}{ l r }
				姓\hfill 名：                   & \underline{{\makebox[6cm][c]{\Author}}}        \\
				学\hfill 号：                   & \underline{{\makebox[6cm][c]{\StudentID}}}     \\
				学\hfill 院：                   & \underline{{\makebox[6cm][c]{\Department}}}    \\
				专\hfill 业：                   & \underline{{\makebox[6cm][c]{\Major}}}         \\
				指\hfill 导\hfill 教\hfill 师： & \underline{{\makebox[6cm][c]{\Supervisor}}}    \\
				职\hfill 称：                   & \underline{{\makebox[6cm][c]{\AcademicTitle}}} \\
			\end{tabular}\\
			\vspace{1em}
			\CompleteYear\hspace*{1em}年\hspace*{1em}\CompleteMonth\hspace*{1em}月
		\end{center}
	\end{textblock}
	
\end{titlepage} %插入内封面
\ClearPageStyle

\restoregeometry
%生成目录
\addtocontents{toc}{\protect\thispagestyle{empty}}
\begin{spacing}{1}
    \tableofcontents
\end{spacing}
\ClearPageStyle
\pagenumbering{Roman}
\thispagestyle{fancy}
\newcommand{\TitleCHS}{ 中文名字 } %中文标题

\newcommand{\TitleENG}{ English name } %英文标题

\newcommand{\Author}{xxx} %作者名字

\newcommand{\StudentID}{xxx} %学号

\newcommand{\Department}{数据科学与工程学院} %学院

\newcommand{\Major}{数据科学与大数据技术} %专业

\newcommand{\Supervisor}{xxx} %导师名字

\newcommand{\AcademicTitle}{xxx} %导师职称

\newcommand{\CompleteYear}{2021} %毕业年份

\newcommand{\CompleteMonth}{7} %毕业月份

\newcommand{\KeywordsCHS}{ 深度学习，xxx，xxx，xxx，分布式训练 } %中文关键词

\newcommand{\KeywordsENG}{ Deep Learning Framework, xxx, xxx, xxx, Distributed Training } %英文关键词

\renewcommand\abstractname{\sffamily\zihao{-3} \TitleCHS}
\phantomsection
\begin{abstract}
	\addcontentsline{toc}{section}{摘要}
	\zihao{5}\rmfamily
	\vspace{\baselineskip} 

	\noindent{\bfseries \rmfamily\zihao{5} 摘要：}

	\par 这里是中文摘要。这里是中文摘要。这里是中文摘要。这里是中文摘要。这里是中文摘要。这里是中文摘要。这里是中文摘要。这里是中文摘要。这里是中文摘要。这里是中文摘要。
	\newline
	\newline
	{\bfseries \rmfamily\zihao{5} 关键词：} \zihao{5}{\rmfamily \KeywordsCHS}
\end{abstract} %生成中英文摘要及关键词
\ClearPageStyle

\thispagestyle{fancy}
\renewcommand\abstractname{\zihao{-3} \TitleENG}
\phantomsection
\begin{abstract}
    \addcontentsline{toc}{section}{ABSTRACT}
    \zihao{5}
    \vspace{\baselineskip}

    \noindent{\bfseries \zihao{5} Abstract:}

    \par English Abstract
    \newline
    \newline
    {\bfseries \zihao{5} Keywords：} {\zihao{5} \KeywordsENG}
\end{abstract} %生成中英文摘要及关键词
\ClearPageStyle
\pagenumbering{arabic}
\pagestyle{fancy} %开始使用页眉页脚
\setcounter{page}{1} %论文页码从正文开始记数

\section{绪论}
 
\subsection{背景}

背景\cite{lecun2015deep}     %正文第一章
\section{SectionB}
        %正文第二章
\section{SectionC}
      %正文第三章
\section{SectionD}

      %正文第四章
\section{SectionE}


             %正文第五章
\section{SectionF}



% SectionFConclusion
\theendnotes %尾注（若没有尾注请将本行删除）
\ClearPageStyle

%生成参考文献
\phantomsection
\addcontentsline{toc}{section}{参考文献}
\printbibliography[title={\centerline{\bfseries\sffamily \zihao {-3}参考文献}}]
\ClearPageStyle

%生成附录
\phantomsection
\addtocontents{toc}{\setcounter{tocdepth}{1}}
\addcontentsline{toc}{section}{附录}
\setcounter{subsection}{0}
\ctexset { subsection = { name={,},number={\arabic{subsection}},format={\rmfamily \zihao {5}} } }
\ctexset { subparagraph = { name={(,)},number={\arabic{subparagraph}},format={\rmfamily \zihao {5}},indent=2em } }
\section*{\centerline{\bfseries \sffamily \zihao{-3} 附录}}

\subsection{实验数据}

123

\subsection{调查结果}

123
\ClearPageStyle

\makeacknowledgement %生成感谢

\end{document} 
